\documentclass[10pt,a4paper]{amsart}
\usepackage[margin=19mm]{geometry}
\usepackage{amsmath,amssymb,mathtools}
\usepackage[scr=rsfs,cal=euler]{mathalfa}
\usepackage{xfrac}
\usepackage[unicode,hidelinks,bookmarks=false]{hyperref}
\newcommand{\ie}{i.e.\ }
\newcommand{\cf}{cf.\ }
\newcommand{\resp}{respectively\ }
\newcommand{\Subsection}{\subsection}
\providecommand{\nobreakdash}{\nobreak\hbox{-}}
\setlength{\emergencystretch}{2em}
\multlinegap0pt
\usepackage{bm}
\usepackage{bbm}
\usepackage{dsfont}
\usepackage{xspace}
\usepackage{cancel}
\usepackage{mathtools}
\usepackage{amsthm}
\makeatletter
\@ifundefined{lemma}{\newtheorem{lemma}{Lemma}}{}
\@ifundefined{proposition}{\newtheorem{proposition}{Proposition}}{}
\makeatother
\DeclareMathOperator*{\argmin}{arg\,min}                   % arg min
\newcommand{\bfA}{{\bf A}}
\newcommand{\bfC}{{\bf C}}
\newcommand{\bfH}{{\bf H}}
\newcommand{\bfL}{{\bf L}}
\newcommand{\bfU}{{\bf U}}
\newcommand{\bfW}{{\bf W}}
\newcommand{\bfZ}{{\bf Z}}
\newcommand{\bfb}{{\bf b}}
\newcommand{\bfe}{{\bf e}}
\newcommand{\bfx}{{\bf x}}
\newcommand{\bfy}{{\bf y}}
\title[Flexible inner-product free Krylov methods for inverse probl]{Flexible inner-product free Krylov methods for inverse problems}
\author{Malena Sabat\'e Landman}
\date{Source: arXiv:2510.18853v1 (CC BY 4.0)}
\begin{document}
\maketitle

\noindent\emph{Source and licence.} Flexible inner-product free Krylov methods for inverse problems. Version arXiv:2510.18853v1 (CC BY 4.0), distributed under the Creative Commons Attribution 4.0 International licence, as recorded in the arXiv licence metadata for this exact version and shown on the abstract page https://arxiv.org/abs/2510.18853v1. Full rights notice: licenses/arxiv-2510.18853-CC-BY-4.0.txt.\par\smallskip

\noindent\textit{Editorial scope.} Counted unit: the Proposition whose source label is \texttt{prop3} in the article (source0.tex lines 760--765, source label \texttt{prop3}), delivered together with the article's own complete original proof of it (source0.tex lines 767--787, one \texttt{proof} environment of 21 lines). No mathematics is written, repaired or re-derived: statement and proof are the article's own text line for line. Required context retained uncounted: eq:functional (lines 459--461); eq:both\_factorization (lines 596--598); eq:IRW\_H\_projection\_approx (lines 626--628); lemma1 (lines 677--685). Every definition, display and label of those lines is retained unchanged. Omitted from this package and not redistributed: 1--458, 462--595, 599--625, 629--676, 686--759, 766--766, 788--1068. Nothing inside a retained excerpt is omitted or altered: every retained body is delivered exactly as the article's lines are written, so no omission receipt of any kind accompanies this package. Citations: the retained excerpts carry no citation command at all, so no bibliography entry of the article is transcribed and no reference is redistributed. Reference adaptation: none was needed and none was made; every reference inside the delivered unit resolves inside it, so no unresolved reference or citation is produced. Printed slips kept literal: no printed word, symbol, hypothesis or number of the counted statement or of its proof was altered. Layout and class: the article's own class is \texttt{article} with packages including geometry, graphicx, amsmath, amsthm, xcolor, algorithm, algpseudocode, amsfonts, caption, subcaption, lipsum, authblk, url, enumitem; the wrapper is the lane's pinned \texttt{amsart} editorial wrapper at 10pt on A4, which restates the theorem-like environments the retained text needs and adds the article's own macro definitions that the retained text needs (\texttt{\textbackslash argmin}, \texttt{\textbackslash bfA}, \texttt{\textbackslash bfC}, \texttt{\textbackslash bfH}, \texttt{\textbackslash bfL}, \texttt{\textbackslash bfU}, \texttt{\textbackslash bfW}, \texttt{\textbackslash bfZ}, \texttt{\textbackslash bfb}, \texttt{\textbackslash bfe}, \texttt{\textbackslash bfx}, \texttt{\textbackslash bfy}), with extra wrapper packages (bm, bbm, dsfont, xspace, cancel, mathtools). The delivered numbering is the unit's own. Assets: no retained span contains a figure environment, a graphics-inclusion command, a PDF-page inclusion, a table or a TikZ picture; no image file is copied into this package, the package declares no asset dependency and no image file is redistributed with it. Licence: version arXiv:2510.18853v1 of this article is distributed under the Creative Commons Attribution 4.0 International licence, as recorded in the arXiv licence metadata for this exact version and shown on the abstract page https://arxiv.org/abs/2510.18853v1. A full rights notice is delivered at licenses/arxiv-2510.18853-CC-BY-4.0.txt. Characters: every retained excerpt is pure ASCII, so the delivered engine, XeLaTeX with its default Latin Modern text font, reports no missing character.
\begin{equation}\label{eq:functional}
\displaystyle \min_\bfx \underbrace{\left\{\| \bfA \bfx - \bfb \|_2^2 + \lambda \|\bfW(\bfL \bfx) \bfL \bfx \|_2^2\right\}}_{f(\bfx)},
\end{equation}

\begin{equation}\label{eq:both_factorization}
    \bfA \, \bfZ_{k} = \bfU_{k+1}  \, \bfH_{k+1,k},
\end{equation}

\begin{equation}\label{eq:IRW_H_projection_approx}
   \min_{\bfy\in\mathbb{R}^k} \| \bfH_{k+1,k}\bfy -\beta \bfe_1 \|_2^2+\lambda\|\bfU^{LU}_k\bfy\|_2^2, \quad \text{for} \quad \bfL^{LU}_k\bfU^{LU}_k = \bfW_k \bfL\,\bfZ_k,
\end{equation}

\begin{lemma}\label{lemma1} Given two functions $h(\bfx)$ and $\hat h(\bfx)$ such that
\[
k_2 h(\bfx) \leq \hat h(\bfx) \leq k_1 h(\bfx), \quad k_1 > 0, \  k_2 > 0,
\]
and define $ \bfx_k=\argmin_{\bfx\in\mathcal{R}(\bfZ_k)} \bar h (\bfx)$. Then, for any $\bfx_{k-1}\in \mathcal{R}(\bfZ_{k-1}) \subset \mathcal{R}(\bfZ_k)$,
\[
\frac{\hat h(\bfx_{k-1})-\hat h(\bfx_{k})}{\hat h(\bfx_{k})} \geq\frac{k_1-k_2}{k_2} \quad  \Rightarrow \quad  h(\bfx_{k-1}) \geq h(\bfx_{k}) 
\]
\end{lemma}

\begin{proposition}\label{prop3} Given a sequence of solutions $\{\bfx_k\}_{k=1,..	}$ obtained with IRW-FCMRH or IRW-FLSLU, then,
for $f(\bfx)$ in \eqref{eq:functional} with a fixed regularization parameter $\lambda$ %if 
\[\frac{\hat q_k(\bfx_{k-1})-\hat q_k(\bfx_{k})}{\hat q_k(\bfx_{k})} \geq \frac{\max(\sigma_1^2(\bfU_{k+1}^{\dagger}),\sigma_1^2((\bfU_{k}^{LU})^{\dagger}))}{\min(\sigma_{k+1}^2(\bfU_{k+1}^{\dagger}),\sigma_k^2((\bfU_{k}^{LU})^{\dagger}))}-1 = \kappa(\bfC)^2-1 \Rightarrow f(\bfx_{k-1}) \geq f(\bfx_k) \]
where $\hat q_k(\bfx_{k-1})$ is the functional minimized at each iteration in the corresponding Krylov subspace, where 
$\bfU_{k+1}$ and $\bfU^{LU}_{k}$ are defined in  \eqref{eq:both_factorization} and \eqref{eq:IRW_H_projection_approx}, respectively, and $\bfC$ is a block diagonal matrix with $\bfU_{k+1}^{\dagger}$ and $(\bfU^{LU}_{k})^{\dagger}$ in its diagonal. 
\end{proposition}

\begin{proof}
Consider the quadratic tangent majorant  $q_k(\bfx)$ of $f(\bfx)$ at $\bfx_{k-1},$ 
\begin{equation*}
   q_k(\bfx) = \| \bfA \bfx - \bfb \|_2^2 + \lambda \|\bfW_k \bfL \bfx \|_2^2 + c_k, \quad \bfW_k =\bfW(\bfL \bfx_{k-1}),
\end{equation*}
where $c_k$ is a constant that does not depend on the solution and $\lambda$ has absorbed possible multiplicative factors. At each iteration of IRW-FCMRH or IRW-FLSLU, one minimizes \eqref{eq:IRW_H_projection_approx}, which corresponds to minimizing 
\begin{equation*}
   \hat q_k(\bfx) = \| \bfU_{k+1}^{\dagger} (\bfA \bfx - \bfb) \|_2^2 + \lambda \|(\bfU_{k}^{LU})^{\dagger} \bfW_k \bfL \bfx \|_2^2 + c_k, \quad \bfW_k =\bfW(\bfL \bfx_{k-1}),
\end{equation*}
for $\bfU_{k+1}$ defined in  \eqref{eq:both_factorization} and $\bfU^{LU}_{k}$ defined in \eqref{eq:IRW_H_projection_approx}. Then,
\begin{eqnarray*}
   \hat q_k(\bfx) \leq \sigma_1^2(\bfU_{k+1}^{\dagger})\, \| (\bfA \bfx - \bfb) \|_2^2 + \lambda \,\sigma_1^2((\bfU_{k}^{LU})^{\dagger} )\, \|\bfW_k \bfL \bfx \|_2^2  + c_k  \leq \underbrace{\max(\sigma_1^2(\bfU_{k+1}^{\dagger}),\sigma_1^2((\bfU_{k}^{LU})^{\dagger}))}_{k_1 \geq 0 } \, q_k(\bfx), \\
 \hat q_k(\bfx) \geq \sigma_{k+1}^2(\bfU_{k+1}^{\dagger})\, \| (\bfA \bfx - \bfb) \|_2^2 + \lambda \sigma_{k}^2((\bfU_{k}^{LU})^{\dagger}) \|\bfW_k \bfL \bfx \|_2^2  + c_k  \geq \underbrace{\min(\sigma_{k+1}^2(\bfU_{k+1}^{\dagger}),\sigma_k^2((\bfU_{k}^{LU})^{\dagger}))}_{k_2 \geq 0 } \, q_k(\bfx).
\end{eqnarray*}
Using Lemma \ref{lemma1} with $h(\bfx) = q_{k}(\bfx) $ and $\hat h(\bfx) = \hat q_{k}(\bfx) $, we have that:
\[\frac{\Delta_q}{\hat q_k(\bfx_{k})} \geq \frac{k_1-k_2}{k_2}= \frac{\max(\sigma_1^2(\bfU_{k+1}^{\dagger}),\sigma_1^2((\bfU_{k}^{LU})^{\dagger}))}{\min(\sigma_{k+1}^2(\bfU_{k+1}^{\dagger}),\sigma_k^2((\bfU_{k}^{LU})^{\dagger}))}-1 \Rightarrow  q_k(\bfx_{k-1}) - q_k(\bfx_{k}) \geq 0.\]
If this condition is met, then a descent step in the original function is also guaranteed:
\[f(\bfx_{k}) \leq q_k(\bfx_{k}) \leq q_k(\bfx_{k-1}) = f(\bfx_{k-1}), 
\]
where the first inequality holds since $ q_k(\bfx)$ is a majorant of $f(\bfx)$ and the last equality holds given the definition of $ q_k(\bfx)$. 
\end{proof}

\end{document}
